\documentclass{article}

\usepackage{graphics} % Required for inserting images
\usepackage{amsmath} % Required for math symbols
\usepackage{amsfonts} % Required for math symbols
\usepackage{color} % Required for \todo notes
\usepackage{svg} % Required for including svg images
\usepackage{setspace} % Required for double space
\usepackage{lineno} % Required for line numbers
\usepackage{hyperref} % Required for figure and eq hyperlinks
%\linenumbers
%\doublespace
\usepackage[margin=1in]{geometry}
\usepackage{authblk}
\usepackage{bm}

\newcommand{\todo}[1]{\textcolor{red}{[#1]}}
\newcommand{\question}[1]{\textcolor{orange}{Q: [#1]}}
\newcommand{\note}[1]{\textcolor{blue}{NOTE: #1}}
\renewcommand{\vec}{\mathbf}
\newcommand{\mat}{\mathbf}
\newcommand{\tensor}{\mathcal}
\newcommand{\new}[1]{\textcolor{purple}{[#1]}}
\newcommand{\eset}{\mathcal{E}}

\title{Visual compressed sensing}

\author[1,*]{Suhyun Ban}
\author[1,*]{Ruchi Gupta}
\author[1,*]{Veronica Hahn}
\author[1,*]{Angus Read}
\author[1,**]{Kameron Decker Harris}
\affil[1]{Western Washington University, Department of Computer Science}
\affil[*]{Alphabetical order}
\affil[**]{Corresponding author: \href{mailto:kameron.harris@wwu.edu}{kameron.harris@wwu.edu}}

\date{\today}


\begin{document}

\maketitle

\begin{abstract}
  The goal of compressed sensing (CS) is to estimate a sparse signal using less samples than its intrinsic dimensionality.
  A large body of existing work shows that random projections are well-suited for this purpose, but these random projections are typically taken to be unstructured.
  Receptive fields in mammalian primary visual cortex (V1) have many properties similar to spatially structured random projections.
  In this paper, we apply these V1-like random projections to image sensing tasks.
  We find that V1-like random projections consistently outperform classical random projection methods. 
  We then look for properties of the design matrix created using V1-like sampling that might explain why this is.
  \todo{why? RIP, etc}
  \todo{wrap-up}
\end{abstract}


\section{Introduction}


Naturally, taking more measurements of a signal gives more data and will result in a better ability to approximate the signal.
The goal of compressed sensing (CS) is to accurately reconstruct an underlying signal using fewer measurements than traditional sampling methods \cite{foucartMathematicalIntroductionCompressive2013, raniSystematicReviewCompressive2018}.
CS is used in several modern domains such as medical and scientific imaging.
Better sensing approaches may translate to decreasing the number of scans needed for generating an accurate picture which reduces exposure to radiation during an X-ray \cite{graffCompressiveSensingMedical2015}.
\todo{bring in brain connections here}

The modern formulation of CS originates from work done by Cand\'es, Romberg, Tao and Donoho in the early 2000's \cite{candesRobustUncertaintyPrinciples2004a, candesOptimalSignalRecovery2006, donohoCompressedSensing2006}.
Cand\'es, Romberg and Tao show that a signal can be exactly reconstructed using $\ell_1$ minimization from an incomplete set of observations given that the underlying signal is formed from a sparse basis, forming the basis of the LASSO regression algorithm.
A related field is sparse coding, which is a technique that looks for a sparse solution in a given dictionary---a set of vectors that may be larger than the dimension of the signal itself.
In the context of imaging, CS seeks to build the image, whereas sparse coding looks for its sparse representation in a wavelet or similar dictionary.

Most natural images do not have a sparse encoding in the pixel basis, as this would require many of the pixels to be black.
Features such as edges and smooth regions mean that they are approximately sparse in an appropriate dictionary \cite{mallat}.
%If this is the case, we know that it is possible to generate an accurate estimate of the image with a limited set of observations \cite{candesRobustUncertaintyPrinciples2004a, candesOptimalSignalRecovery2006,donohoCompressedSensing2006}.
JPEG uses the discrete cosine transform (DCT) for lossy image compression by encoding $8 \times 8$ non-overlapping windows and using a sinusoidal Fourier basis within each window.
After the DCT coefficients are extracted, the top coefficients retained and quantized to optimize storage space.
When the image is decompressed, the inverse of this process is applied to generate a good approximation of the original image \cite{wallaceJPEGStillPicture1992}.
Another method of generating a basis for continuous functions from discrete observations is discrete wavelet transform (DWT).
Wavelets allow for localization in space as well as frequency;
JPEG achieves spatial localization with windowing.
Both DCT and DWT offer bases for sparse image encoding.

In animals, visual neurons that receive input from the retina respond selectively to certain features.
We characterize these preferred features in terms of their receptive fields.
The receptive fields of cells in primary visual cortex (V1) are preferred patterns---usually edge-like---within localized spatial regions \cite{hubelReceptiveFieldsBinocular1962}.
Mathematical models of V1 typically treat the receptive fields as weights.
Olshausen and Field \cite{olshausenSparseCodingOvercomplete1997} found that the receptive fields learned by applying the sparse coding algorithm to natural images appear very similar to receptive fields measured from the V1 cortex.
These results indicate that if the V1 represents visual data via an over-complete dictionary, we will reasonably be able to approximate the original signal.
These receptive fields have been refined and improved over millions of years of natural selection.
While it's debatable whether brains actually perform a something akin to CS, it is logical to think that we can take inspiration from these systems and try to implement them in a non-natural context.
Hopefully, this will let us both improve upon existing methods of sampling as well as gain some more insight into how these complex biological systems function.

Complementary to \cite{olshausenSparseCodingOvercomplete1997},
\cite{pandeyStructuredRandomReceptive} showed that mouse V1 receptive fields were well-characterized by a random model.
This model draws from a Gaussian process (GP) with a simple covariance function that leads to localized and smooth receptive fields.
With only 2 tunable parameters, the GP model captured much of the covariance structure of the V1 data and led to good performance in visual classification tasks.

%\todo{say "sparse coding" and talk about the shape of receptive fields, like wavelets}.

%\todo{summarize what did for V1 random features and our main results, even the ones that are for supervised learning rather than just CS}

%\note{sources: Pandey et al paper and find right references. Chris Rozell is one to include.}

In this paper, we will compare the effectiveness of V1-like structured random observations against unstructured random sampling.
Measurements are generated by taking images in the pixel basis and applying a randomly-generated sensing matrix.
These observations are then linearly transformed in the hope that the image data has a sparse representation in this new basis.
The most significant basis functions are selected using LASSO and are then transformed back into the pixel-basis, providing an estimation for the original data. 
We test the effect on error of forcing sparsity in the new basis, as well as the impact of varying the number of observations generated.
We compare the mutual coherence of the V1-like design matrix against the Gaussian and pixel matrices. This is a value in matrix theory that measures the maximum similarity between any two columns of a matrix, and low MC values are provably better for compressed sensing \todo{cite classics}. 
We found that the V1-like design matrices consistently had higher mutual coherence.
%We suspect the cause of the high MC value is the ``blobby" nature of the observations.\todo{better way to say this?} 
We then examine the singular value decomposition (SVD) of the design matrix.
Using the singular values, we look at the principal image components of the reconstructions in the right singular vector basis of the design matrix. 
These image components give us insight into which coefficients are determined to be most relevant to reconstruction and how they behave.
More research is needed as to why, but our results indicate that generating measurements that are inspired by structured receptive fields in V1 is more effective in the context of compressed sensing than other methods of modeling unstructured random receptive fields.


\section{Methods}

Now we describe the image data that we use for our experiments, the competing methods of random projection, algorithm for CS in DCT and DWT bases, the specific experiments that were performed, and an investigation into properties of design matrices using V1-like measurements. \new{new}


\subsection{Images}

We use four $512 \times 512$ test images from \url{https://sipi.usc.edu/database/database.php?volume=misc} and \url{https://ieeexplore.ieee.org/document/258085} \cite{shapiro} for our experiments: \textit{Baboon}, \textit{Barbara}, \textit{Boat}, and \textit{Peppers}.
These are commonly used in CS papers for image reconstruction \cite{parkaleApplicationCompressedSensing2021, sermwuthisarnFastImageRecovery2009}.
Two of the images are black and white (\textit{Barbara} and \textit{Boat}), and two are color (\textit{Baboon} and \textit{Fruits}).
They also exhibit a variety of patterns and textures such as the stripes of various size and orientation in the shawl of \textit{Barbara}, the mesh grid in \textit{Fruits} as well as the cables and rigging in \textit{Boat}.
We can then analyze how effective the sampling method is for estimating signals of varying frequency.
We use $\vec x \in \mathbb{R}^d$ to refer to the flattened $d_1 \times d_2$ image or patch that we are reconstructing, with $d = d_1 d_2$.
\todo{verify following:}
When the image has multiple color channels, we sample each color channel independently.


\subsection{Measurement Methods}

We generate a random sensing matrix $\mat W \in \mathbb{R}^{n \times d}$
using either pixel, Gaussian white noise, and V1-like methods.
The $n$ rows in $\mat W$ each correspond to a measurement of the image.
The sensing matrix $\mat W$ contains the weights which we apply to the true pixel data to generate the vector of measurements
\begin{equation}
  \vec{y} = \mat W \vec{x}, \quad \vec y \in \mathbb{R}^n.
  \label{eq:weighted_observations}
\end{equation}



\subsubsection{Pixel Measurements}

Pixel sampling is an unstructured method that returns the value of a random subset of pixels.
The $n \times d$ weight matrix $\mat W$ is formed by first initializing all entries to be 0.
For each row, we select a random pixel from the $d$ possibilities with replacement and set the corresponding matrix entry equal to $\sqrt{d}$.
The scaling is chosen so that each row $\vec w_i$ has norm $\|\vec w_i \|_2^2 = d$ to be consistent with the other sampling methods.
%In both methods, we perform Gaussian sampling and use a matrix $\mat I$ or $\mat C$ to represent covariance between pixels.
% Both matrices have the property tr($\mat I$) = tr($\mat C$) = $d$ which then means that each row $||\vec w_i||^2_2 = \sqrt{d}$.
% We simply weight a single pixel in each row of $\mat W$ by $\sqrt{d}$ to force the same property that $||\vec w_i ||^2_2 = \sqrt{d}$.
% The variable $a$ is used in Equation \hyperref[eq:v1_covariance]{\ref{eq:v1_covariance}} to
% scale $\mat C$ so the condition holds that tr($\mat C$) = $d$.


\subsubsection{Gaussian White Noise Measurements}

Gaussian white noise sampling also generates unstructured random receptive fields.
This method samples each row of weight matrix $\mat W = [ \vec w_1, \ldots, \vec w_n ]^T$ from an isotropic normal distribution
%\begin{equation}
$\vec w_i \sim N(0, I_d) $.
%  \label{eq:whitenoise} 
%\end{equation}
Since different spatial locations are uncorrelated, we call this an unstructured random receptive field.
Note that $\mathrm{Tr}(\mat I_d) = d = \mathbb{E}[ \| \vec w_i \|_2^2 ]$.



\subsubsection{V1-like Gaussian Measurements}

The V1-like sampling method uses correlated Gaussian vectors to generate structured receptive fields.
This random model was shown to capture many properties of neurons in V1 \cite{pandeyStructuredRandomReceptive}.
In this method, the rows of $\mat W$ are drawn from
$  \vec w_i \sim N(0, \mat C) $,
with pixel covariance matrix $\mat C \in \mathbb{R}^{d \times d}$.
The covariance matrix is generated by sampling the covariance function
\begin{equation}
  C(\vec t, \vec {t'}) =
  a \,
  %\cdot
  \exp \left(-\frac{\|\vec t-\vec t'\|^2}{2f^2} \right)
  %\cdot
  \exp \left( -\frac{\|\vec t - \vec c\|^2 + \|\vec t' - \vec c\|^2}{2s^2} \right)
  \label{eq:v1_covariance}
\end{equation}
on the $d_1 \times d_2$ pixel grid.
The parameters
$\vec t \in \mathbb{R}^2$ and
$\vec t' \in \mathbb{R}^2$ are the 2-d locations of pixels in the image, and
$\vec c \in \mathbb{R}^2$ is the center of the receptive field.
The center $\vec c$ is drawn uniformly at random from the $d_1 \times d_2$ grid of pixels.
%We set, $\vec c = (\lfloor d_1 / 2 \rfloor, \lfloor d_2 / 2 \rfloor)$ to the center of the window \todo{remove floors?}.
The scale parameter $a$ is chosen so that $\mathrm{Tr}(\mat C) = d$.
The spatial bandwidth $f$ allows us to control the length scale of variation of the receptive fields; increasing $f$ leads to smoother within the receptive fields.
Field size $s$ sets the spatial extent of non-zero receptive field weights; increasing $s$ leads to a larger receptive field.
We treat $s$ and $f$ as hyperparameters of the sampling method and optimize them in experiments.

%\todo{The only extra thing to add here is that in Pandey et al we showed that this form of covariance is diagonalized by Hermite wavelets, and the eigenvalues can be thought of as performing low-pass filtering. Rephrase that your own way}

\subsection{Sparse Coding Bases}

We consider two possible representations in which the image is assumed to be sparse: DCT or DWT.
with $\vec z \in \mathbb{R}^p$ a sparse vector.
Our observations $\vec y = \mat W \vec x$, so the observation process from the sparse representation depends on a matrix
$\mat A = \mat W \mat U$, $\mat A \in \mathbb{C}^{n \times p}$.
By our convention, the matrix $\mat U^T$---acting from the left on a column vector---is the forward transform and $\mat U$ is the inverse transform with $\mat U \in \mathbb{C}^{d \times p}$ its matrix representation.
Assuming the image $\vec x$ is sparse in this basis is equivalent to
$\vec x \approx \mat U \vec z$.
We write $\vec z^* = \mat U^T \vec x$ so that $\vec y = \mat A \vec z^*$.

To compute $\mat A$, we avoid building the matrix $\mat U$ and instead leverage fast transform techniques:
We reshape the sensing matrix $\mat W$ into an $n \times d_1 \times d_2$ tensor, operate on the last two dimensions with 2-d forward transforms, then reshape to get $\mat A$.

The DCT is unitary, so $\mat U$ is always square with $p = d$.
The DWT on the other hand uses an overcomplete dictionary and may have $p \geq d$.
DWT differs from DCT in that it sums over generated wavelet functions rather than cosines.
Wavelets are localized in both space and frequency domains, unlike cosine waves which repeat infinitely and only retain frequency.
The DWT implementation we use \cite{dwtpak} introduces the  \texttt{level} hyperparameter.
At each level, the coefficients are split up into higher frequency and lower frequency coefficients.
This means that as the number of levels increases, more splits occur and the attention to individual detail increases.


\subsection{Reconstruction}
Once we have formed $\mat A \in \mathbb{R}^{n \times p}$, we solve for a sparse encoding $\vec z \in \mathbb{R}^p$ of our observations $\vec y$.
A full description of the LASSO formulation used in some of our experiments, along with its KKT optimality conditions, is given in Appendix~\ref{app:additional-experiments}.
The entries in $\vec z$ correspond to the most important cosine or wavelet functions for reconstructing our weighted observations $\vec y$.
To reconstruct the image, we reshape the coefficients $\vec z$ to the appropriate shape ($d_1 \times d_2$ array for DCT or using an unravel command for DWT) and apply the inverse transform to get the estimate of the image $\hat{\vec x} = \mat U \vec z$.

\subsection{Windowing for Large Images}
Until now, we have been talking about reconstructing the entire image at once.
In practice, it is more efficient and accurate to reconstruct the image in patches.
We do this by sliding a $d_1 \times d_2$ non-overlapping window across each color channel of the image.
We then generate weighted observations for those pixels and attempt to reconstruct the patch.
For all data in this paper, a new sensing matrix $\mat W$ is computed after each patch to model a different group of receptive fields for sensing the next patch.
In this way, each patch has a unique sensing matrix $\mat W$ and as such, has unique receptive fields.
We also experimented with using a fixed matrix $\mat W$.
This yielded similar results when compared to generating a new $\mat W$ for each patch, though these results are not shown.


\subsection{Experiments}
We tested the effectiveness of these different random sampling methods by performing a grid search over many parameters.
The results were then analyzed to see how various parameters effect the accuracy of reconstruction.
The hyperparameters which we vary are the LASSO penalty $\alpha$, field size $s$, spatial bandwidth $f$, and \texttt{level},
We also test the impact of varying $n$, and the patch dimensions $d_1$ and $d_2$.
In all sampling scenarios and trials we use $d_1 = d_2$ for square patches.
All DCT trials use the same hyperparameters: $s \in \{50, 100, 200\}$ and $f \in \{2, 4, 6, 8\}$, $\alpha \in \{0.01, 0.1, 1, 10\}$, $d \in \{8 \times 8, 16 \times 16, 32 \times 32\}$.
In order to make reconstructing different patch sizes comparable, we used the same sampling percentage per patch, but this results in different $n$ for different patch sizes.
Our sampling percentage were $12.5\%, 21.9\%, 31.3\%, 40.6\%, 50\%$.
For $d = 64, 256, 1024$ this resulted in $n \in \{8, 14, 20, 26, 32\}, \{32, 56, 80, 104, 128\},$ and $\{128, 224, 320, 416, 512\}$ respectively. 
DWT includes 2 unique parameters: ``db2'' Daubechies wavelets, and level $ \in \{1, 2, 3\}$.
We then generate our search space by computing all possible parameters between those sets, and repeating each combination 10 times.
The accuracy of each construction is quantified by computing the root mean square error: $\mathrm{Error} = \frac{1}{\sqrt{d}}||\vec x - \mat U \vec z||_F$

\subsection{V1-like Design Matrix Experiment Methods}

\subsubsection{Mutual Coherence}

Let $M$ be the mutual coherence of the design matrix $\mat{A}$.
Let $\vec{a}_1, \ldots, \vec{a}_{d}$ be the column vectors of $\mat{A}$. 
Then,
\begin{equation}
    M = \max_{i \neq j} \frac{|\vec{a}_i^T \vec{a}_j|}{\|\vec{a}_i \|_2 \|\vec{a}_j \|_2}.
\end{equation}
Equivalently, we can say that mutual coherence is the maximum column dot product of $\mat{A}$ with itself ($\mat{AA}^T$) after the column vectors have been normalized. 
We then show the histograms of the dot products, the tails of which are the largest values, for design matrices based on V1-like, Gaussian, and pixel sampling of the same patch.


\subsubsection{Singular Value Decomposition and principal components}

Let $\mat{A} = \mat{L} \bm{\Sigma} \mat{R}^T$
be the {\em full} SVD of the observation matrix, 
where the $n \times n$ matrix $\mat{L}$ and 
$p \times p$ matrix $\mat{R}$ orthonormal left and right singular matrices.
Since the observations are multivariate Gaussian, there are $n$ nonzero singular values and thus $p -n $ zero singular values.

After reconstructing using $\mat{A}$, we can now define the estimated sparse encoding vector $\vec{z}$ in terms of the SVD of A. 
\begin{equation}
\vec{z} = \sum p_i \vec{r}_i = \mat{R} \vec{p}, \quad \mat{R} = [ \vec{r}_1, \ldots, \vec{r}_d], 
\end{equation}
Then rotate \textbf{z} into the $R$ basis.
 \begin{equation}
     \vec{p} = \mat{R}^T \vec{z}.
 \end{equation}
We call the elements of vector $\vec{p}$ the true image components in the PC basis of $\mat{A}$ and $\hat{\vec{p}}$ the estimated image components.
We study the principal components by showing them in a scatter plot as well as comparing the estimated principal components and the true ones from patch 58 of \textit{Barbara}.


\subsection{Software}
We used the code produced by \cite{pandeyStructuredRandomReceptive} to generate structured V1 weights: \url{https://github.com/BruntonUWBio/structured-random-features}.
We also significantly used numpy \cite{harrisArrayProgrammingNumPy2020} pywavelets \cite{leePyWaveletsPythonPackage2019}, scikit-learn \cite{pedregosaScikitlearnMachineLearning} and scipy \cite{virtanenSciPyFundamentalAlgorithms2020} to aid in basis-change computation and reconstruction of images.


\section{Results}

\begin{figure}[h!]
  \includegraphics[width=\textwidth]{images/zoomed.jpg}
  \centering
  
  \caption{
    Here, we show the zoomed in DCT reconstruction of \textit{Barbara} using each sampling method.
    On the left is the original image and each circle on the right corresponds to a zoomed in DCT reconstruction.
    We used parameters $n = 256$, $\alpha = 1$, $d = 32 \times 32$ which corresponds to a sampling percentage of $25\%$.
    The V1 parameters were $s = 50$ and $f = 6$.
    The yellow ovals show the stains on the shawl and pants, which is only visible in the original image and the V1 reconstruction.
    Patch boundary artifacts are mostly evident in the pixel reconstruction and are highlighted by the blue oval.
    Green ovals are used to draw attention to the stripes on the shawl and each sampling method's effectiveness at reconstructing the pattern.
    \label{fig:zoomed}
  }
\end{figure}

\begin{figure}[h!]
  \includegraphics[width=\textwidth]{images/dct_reconstructions.pdf}
  \centering
  \caption{
    We present a reconstruction of each image in the test set using a DCT transformation, with the original image depicted in the top row. 
    Rows 2-4 are the reconstruction of each image with each row corresponding to different method of generating receptive fields.
    We used parameters $n = 256$, $\alpha = 1$, $d = 32 \times 32$.
    The V1 parameters were $s = 50$ and $f = 6$.
    These reconstructions are generated from a sampling percentage of $25\%$.
    \label{fig:dct_reconstructions}
  }
\end{figure}



% \begin{figure}[h!]
%   \includegraphics[width=\textwidth]{images/dwt_reconstructions.jpg}
%   \centering
  
%   \caption{
%     Each column in this figure is a different image we used to test compression and reconstruction and the top row is the original image.
%     Rows 2-4 are the DWT reconstruction of each image with each row corresponding to different method of generating receptive fields.
%     We used parameters $n = 256$, $\alpha = 1$, $d = 32 \times 32$.
%     The V1 parameters were $s = 50$ and $f = 6$.
%     The DWT parameters were \texttt{level} = 3 with Daubechies ``db2'' wavelets. 
%     These reconstructions are generated from a sampling percentage of $25\%$.
%     \label{fig:dwt_reconstructions}
%   }
% \end{figure}



In our analysis of the results, we test the biologically-inspired V1 sampling method (structured receptive fields) against the pixel sampling and Gaussian white noise methods (unstructured receptive fields).
In Figure \hyperref[fig:zoomed]{\ref{fig:zoomed}} we present the original image on the left, and show a zoomed in version of each reconstruction.
This figure illustrates each sampling method's effectiveness in capturing various patterns and textures.
We can see that in the V1-like sampling, the striped pattern on the shawl and pants gets smoothed out and is not depicted in the reconstruction.
However, the V1 sampling does a better job than both of the unstructured random sampling methods at reconstructing the stain on the pants and the shawl.
We can see that in all of the reconstructed images, there are some notable effects on the boundary of each patch.
Figure \hyperref[fig:dct_Reconstruction's]{\ref{fig:dct_reconstructions}} depicts a reconstruction of each image in our test set with each sampling method using DCT.
In \textit{Baboon} it is hard to see a significant difference between each sampling method, but for V1 the text on the back of \textit{Boat} is more readable, and the rigging is more noticeable.
In \textit{Fruits} the most noticeable differences are the texture of the grapes, as well as the mesh grid in the top left of the image.
Figure \hyperref[fig:dwt_reconstructions]{\ref{fig:dwt_reconstructions}} is the same as the previous figure, except that it uses DWT to change basis rather than DCT.
Clearly, we can see that this gives a significantly worse reconstruction of the image in nearly every case, exhibiting very strange patterns at the boundary of each patch.
The V1-like sampling method  does relatively well but still has some noticeable effects on each patch boundary.

Examining the mutual coherence for various design matrices reveals that V1-like sampling results in a worse MC value  every time. 
The histogram in figure \hyperref[ref:dot_products]{\ref{fig:dot_products}} shows that V1-like sampling has both a worse maximum dot product and more dot products which are higher than the maximum of both Gaussian and pixel measurements. 
The consistency of this property is shown in the same figure as a box plot. 
We move on to studies of the singular values and image components components of $\mat{A}_{58}$.


\subsection{Varying Number of Measurements}

For our analysis in which we vary $n$ we used a fixed patch size of $d = 32 \times 32$.
We then compute the mean reconstruction error averaged over repetitions for each $n$ with a $32 \times 32$ window.
The optimal $\alpha$ is selected and, if applicable, the optimal DWT and V1 hyperparameters level, $s$, $f$ for each $n$. 
Once we have found the  on-average optimal parameters, we can plot the error as a function of $n$ as seen in Figure \ref{fig:vary_n}.
This plot has error bars to show the variation between trials, however the amount of variation is so small that the error bars are nearly too small to see.
For our experiments, we see that varying $n$ results in a better reconstruction for higher $n$.
This is to be expected, as increasing $n$ is increasing the number of observations that we are able to generate per patch.
Therefore, as we gain more information about the true pixel data, we are able to more accurately estimate it.


\begin{figure}[h!]
  \includegraphics[width=\textwidth]{images/varying_n.jpg}
  \centering
  \caption{
    Error versus number of samples $n$ with a fixed patch size of $d = 32 \times 32$.
    The error of reconstruction decreases as we increase the number of measurements that are taken, which is to be expected.
    \label{fig:vary_n}
}
\end{figure}


\subsection{Varying LASSO Penalty}
To analyze the effect of varying the $\alpha$ penalty we fixed $n = 256$ and patch size $d = 32 \times 32$, which is a $25\%$ sampling percentage.
We compute the mean reconstruction error averaged over repetitions for each $\alpha$ with $n = 256$ and a $32 \times 32$ window.
We then select the optimal DWT and V1 hyperparameters $\texttt{level}$, $s$, $f$ for each $\alpha$ if applicable.
Once we have found the optimal hyperparameters, we can plot the error as a function of $\alpha$.
We again include error bars to show variation across trials, which appears to be almost nonexistent.
The V1-like method does much worse than the other two methods with small $\alpha$, but when $\alpha \geq 0.1$, V1 performs better.
We find that each method of sampling has an optimal $\alpha$, typically 1 or 0.1 as seen in Figure \hyperref[fig:alpha]{\ref{fig:alpha}}.
This aligns with results from previous work, as the goal is to gain a sparse representation of a basis and if $\alpha$ is zero, then we are doing simple least-squares regression which does not necessarily give a sparse solution.



% \begin{figure}[h!]
%   \includegraphics[width=\textwidth]{images/varying_alpha.jpg}
%   \centering
%   \caption{
%     We show how the error varies with $\alpha$ with a fixed patch size of $d = 32 \times 32$ and $n = 256$.
%     For very small $\alpha$ we note that V1 does significantly worse than both other methods for DCT and DWT.
%     However, with $\alpha \geq 0.1$ we see that V1 does better than both unstructured sampling methods.
%     These results indicate that the optimal $\alpha$ is either $0.1$ or $1$.
%     \label{fig:alpha}
% }
% \end{figure}

\subsection{Varying Patch Dimensions}

For our analysis of varying the patch dimensions, we used a fixed sampling percentage of $12.5\%$.
We compute the mean reconstruction error averaged over repetitions for each $\alpha$ with $n = 256$ and a $32 \times 32$ window.
We then select the optimal V1 hyperparameters $s$, $f$ for each patch size if applicable.
Once we have found the optimal parameters, we can plot the error as a function of the patch dimensions and include error bars to show variation across trials.
Patch size does not seem to have a large effect on the error of reconstruction, as it seems to be roughly the same for all patch sizes that we tested as seen in Figure \ref{fig:vary_filter_dim}.
However, we found that the smaller filter size does impact the runtime of the reconstruction.
This is evident from Table~\ref{table:times} which shows the average time to reconstruct \textit{Baboon} for varying patch dimensions.


\begin{figure}[h!]
  \includegraphics[width=\textwidth]{images/varying_filter_dim.jpg}
  \centering
  \caption{
    Error variation with patch size for a fixed sampling percentage of $31.25\%$, or $n = 20, 80, 320$ for patch sizes $d = 8 \times 8, 16 \times 16, 32 \times 32$, respectively.
    \label{fig:vary_filter_dim}
  }
\end{figure}

\subsection{Properties of V1-like Gaussian Measurements}

\begin{figure}[t!]
    \centering
    
    \includegraphics[width=0.45\linewidth]{A_images/patch_58/png/DotProductHist_patch_58.png}
    \includesvg[width=0.45\linewidth]{A_images/patch_58/svg/MC_patch_58.svg}
    
    \caption{
    On the left are the column-wise dot products for patch 58, the magnitude of entries in $\mat{A} \mat{A} ^T$, viewed as a histogram.
    On the right is a comparison of several $M$ for each observation type.
    We took 100 different measurements of the same patch and found the mutual coherence for every design matrix, then record them as a box plot.
    Despite V1-like sampling having an excellent showing in other experiments, its design matrices consistently have worse MC than the other types.
    We used parameters $n = 256$, $\alpha = 1$, $d = 32 \times 32$.
    The V1 parameters were $s = 50$ and $f = 6$.
    }
    
    \label{fig:dot_products}
\end{figure}

\subsubsection{Principal components of signals and observation spectra}

\begin{figure}
    \centering
    \includegraphics[width=0.5\linewidth]{A_images/patch_58/png/SVD_patch_58.png}
    
    \caption{
    From the singular value decomposition of $\mat{A}_{58}$, we plotted the elements of $\bm{\Sigma}$ directly. 
    We used parameters $n = 256$, $\alpha = 1$, $d = 32 \times 32$.
    The V1 parameters were $s = 50$ and $f = 6$.
    }
    \label{fig:SVD_spectra_58}
\end{figure}

\begin{figure}
    \centering

   \includegraphics[width=\linewidth]{A_images/patch_58/png/pc_per_method_patch_58.png}

   \includegraphics[width=\linewidth]{A_images/patch_58/png/PC_scatter_patch_58.png}
   
    \caption{
    Above is a scatter plot of the elements of $\hat{\vec{p}}_{58}$ of \textit{Barbara}. 
    Below that is a comparison between the estimated principal components and the true ones taken from the raw image of \textit{Barbara}.
    The dotted line represents where $y=x$.
    We used parameters $n = 256$, $\alpha = 1$, $d = 32 \times 32$.
    The V1 parameters were $s = 50$ and $f = 6$.}
    \label{fig:pcScatterCompare_58}
\end{figure}

\begin{figure}
    \centering

    \includegraphics[width=\linewidth]{A_images/patch_58/png/smoothed_error_cdf_patch_58.png}
    
    \caption{
    This is the cumulative sum of error between the estimated and true image components for each observation type.
    We used parameters $n = 256$, $\alpha = 1$, $d = 32 \times 32$.
    The V1 parameters were $s = 50$ and $f = 6$.
    }
    \label{fig:PCError_single}
\end{figure}

We plot singular values of each observation type in Fig.~\ref{fig:SVD_spectra_58}. 
The Gaussian values have a Marchenko-Pastur distribution \cite{}, pixel follows a step function that reflects the frequency of observing a given pixel 0, 1, 2, etc.\ times, and V1-like has an exponentially decaying spectrum \cite{pandeyStructuredRandomReceptive}.
Due to these differences, V1-like sensing matrices have lower dimensionality than pixel or Gaussian.

Fig.~\ref{fig:pcScatterCompare_58} shows a scatter plot of the principal components of patch 58 of \textit{Barbara} and a plot of them with respect to the real principal components of \textit{Barbara}. 
While the V1-like principal components follow a similar sloping pattern to their singular values, the Gaussian and pixel are comparatively random.
For V1, Fig.~\ref{fig:pcScatterCompare_58} shows the leading principal components are very close to the true principal components. 
Finally, we present a graph showing the cumulative sum of error between the real and estimated principal components. 
In most cases, V1-like sampling has a smaller cumulative error, like in Fig.~\ref{fig:PCError_single}, however in Appendix A we show two examples where the cumulative sum of the error is worse.

\todo{talk about patches (refer to supp) where it doesn't work/different answers come out}

\todo{Kam adds math}


\section{Discussion}

We find that between random unstructured receptive fields, we do not see much difference in reconstruction effectiveness.
Though they are relatively similar in performance, the pixel method allows us a sparser observation matrix $\mat W$ than weights from Gaussian distribution.
We also find that generating weights modeled on a random structured receptive field gives a more accurate reconstruction than both methods of unstructured random sampling.
This is evident in our plots, which show that in every scenario, V1 sampling outperforms both benchmark sampling methods.
In all figures, it is also noticeable that DWT always does worse than DCT.
The main issue that can be observed from the DWT reconstructions is the strange patch boundary effects noticeable in Figure {\ref{fig:dwt_reconstructions}}, which we are unsure how to fix.
These results indicates that generating sample weights with covariance between pixels is more effective than generating weights that have no covariance between pixels.
Filter window size does not seem to have an effect on error of reconstruction, but we find that it does have a rather significant effect on the runtime of reconstruction: smaller filter sizes are significantly faster.
Because the accuracy of reconstruction is so similar between different filter sizes, we note that it is much better to use a smaller filter in order to save compute-time as seen in Table{\ref{table:times}}.
\begin{table}[h!]
  \centering
  \begin{tabular}{| c c c |}
    \hline
    Patch Size & $n$ & Average Reconstruction Time (s) \\ 
    \hline\hline
    32x32 & 100 & 25.32 \\  
    \hline
    16x16 & 100 & 24.90\\
    \hline
    8x8 & 100 & 8.69 \\
    \hline
  \end{tabular}
  \caption{
  \todo{Remove table}
  Average reconstruction time for varying patch dimension over 50 trials for \textit{Baboon} with $n=100$ and $\alpha = 1$ using the pixel sampling method.}
  \label{table:times}
 
\end{table}

We also find that the current theory of mutual coherence does not seem to apply to V1-like sensing.
The maximum similarity between column vectors of V1-like design matrices far exceed that of matrices using Gaussian and pixel sampling.
However, because the error of reconstruction is lower than the other observations, there is evidence of some property of V1 sampling that is yet to be found. 
The singular value decomposition of these $\mat{A}$ matrices seem to support this theory.
Since the singular values of those matrices follow the negative sloping pattern found in figure \hyperref[fig:SVD]{\ref{fig:SVD_spectra_58}}, it indicates that far more coefficients\todo{basis functions?} selected during LASSO are close to 0.\todo{This needs a better explanation.}
\todo{something something sparsity is really good it's the whole point}
The principal components from V1-like sampling have a greater error in later ranks because so many more coefficients are zero-ed out, but the leading principal components are more accurate.
V1-like image components have less total cumulative error in most cases.
\todo{reference math}
This tells us that the V1-like coefficients that are relevant to reconstruction are closer to the true coefficients than ones coming from Gaussian and pixel sampling.
Meanwhile, we are able to ignore more signals during reconstruction while still creating a more accurate representation of the image.
\todo{what else to add here?}

Our results show that in almost every scenario V1-like sensing outperforms the other baselines, though the specific properties that make this true are still unknown.
Further research is needed in that area.
Nonetheless, this indicates that the structured random sampling method is more effective at capturing an underlying signal than simply using random sampling.
These findings have some interesting implications in multiple domains.
For artificial neural networks, which often initialize weights randomly, there may be improvements to be made by studying the potential effect of structured weights in the model architecture.
In a medical context, this method of structured random sampling can be used in real-world applications of CS, such as generating observations for MRI scans or x-ray in a more optimized way.
We hope that our findings here may inspire future work into using structured random sampling for CS in these domains, and serve as an improvement on current methods of sampling.
\question{I feel that this conclusion is very weak, do you have any suggestions on how to improve it?}

\bibliographystyle{plain} % We choose the "plain" reference style
\bibliography{refs} % Entries are in the refs.bib file


\newpage
\appendix

\section{Acknowledgements}

\subsection{Author Contributions}
Computer experiments: SB, RG, VH, AR, KDH.
Theory: KDH.
Figure creation: RG, VH, AR, KDH.
Original draft: AR.
Writing: VH, AR, KDH. 
Supervision: KDH.

\subsection{LLM Use Disclosure}

Various models from Claude code were used to clean up handwritten software, implement experiments, analyze results, and generate figures.
These tools were also used to explore the literature when looking for applicable theory to this area.
Computer experiments and mathematical analysis have been checked by the authors, and the authors take full responsibility for the results in this paper.

\section{Additional Experiments}
\label{app:additional-experiments}

\subsection{Reconstruction Using LASSO}
Once we have formed $\mat A \in \mathbb{R}^{n \times p}$, we can find the sparse encoding from our observations using LASSO:
\begin{equation}
\vec z =
\arg \min_{\vec z}
\frac{1}{2n} \| \vec{y} - \mat A \vec z \|_2^2
+ \alpha \| \vec z \|_1.
\end{equation}
Once this regression task is solved, we have a sparse vector $\vec z \in \mathbb{R}^p$.
The entries in $\vec z$ correspond to the most important cosine or wavelet functions for reconstructing our weighted observations $\vec y$.
To reconstruct the image, we reshape the coefficients $\vec z$ to the appropriate shape ($d_1 \times d_2$ array for DCT or using an unravel command for DWT) and apply the inverse transform to get the estimate of the image $\hat{\vec x} = \mat U \vec z$.

\subsection{LASSO KKT Conditions}
Cite Tibshirani paper.
LASSO solution satisfies Karush-Kuhn-Tucker (KKT) conditions:
\begin{align}
    \vec{a}_i^T (\vec{y} - \mat{A} \vec{z})
    &= n \lambda \, \mathrm{sign}(z_i) & \mbox{$z_i \neq 0$}\\
    |\vec{a}_i^T (\vec{y} - \mat{A} \vec{z})|
    &<  n \lambda & \mbox{$z_i = 0$}
\end{align}
If we write $\eset$ for the set of indices with $z_i \neq 1$, set $|\eset|=S$, and define $\vec{s} = (\mathrm{sign}(z_i))$ where $\mathrm{sign}(z) = 0$ for $z=0$, and use the source condition on the observations,
this is equivalent to the matrix system
\begin{align}
    \mat{A}_\eset^T \mat A(\vec{z}^* - \vec{z}) &= n \lambda \vec{s}_\eset \\
    | \mat{A}_{\eset^c}^T \mat A (\vec{z}^* - \vec{z}) | &< n \lambda,
\end{align}
where $\mat{A}_\eset$ is the $n \times S$ matrix with columns from the set $\eset$.
Likewise we get
$\vec p = \mat R^T (\vec z_\eset; \vec 0) = \mat R_\eset^T \vec z_\eset$, where $\mat R_\eset$ contains the rows of $\mat R$ indexed by $\eset$.
Since $\mat A_\eset = \mat U \mat \Sigma \mat R_\eset^T$,
The solution for the nonzero coefficients satisfies:
\begin{align*}
    \vec z_\eset &=
    (\mat A_\eset^T \mat A_\eset)^{-1}
    \left(
    \mat A_\eset^T \mat A \vec z^* - n \lambda \vec s_\eset
    \right)   \\
    &=
    (\mat R_\eset \mat \Sigma \mat R_\eset^T)^{-1}
    \left(
    \mat R_\eset \mat \Sigma^2 \vec p^* - n \lambda \vec s_\eset
    \right)
    &=
\end{align*}

\section{Generality of Various Experiments Across Patches}

\begin{figure}[t]
    \centering
    \includegraphics[width=\linewidth]{A_images/barbara_patches.png}
    \caption{
    These are the $32\times32$ pixel patches we show results for in this paper.
    We believe patch 58 to be a good example because of the complexity of pixel signals found in her face and eye.
    }
    \label{fig:barbara_patches}
\end{figure}

The following figures show the generality of many metrics talked about in this paper.
Some metrics, like the cumulative sum of errors in Fig ~\ref{fig:PCError}, do worse for patches 169 and 233.
However, the majority of patches do well in these experiments.

\begin{figure}[h]
    \centering
   
   \includegraphics[width=\linewidth]{A_images/all_patches/png/MC_all_patches.png}      
   
    \caption{
    We created 100 design matrices for each observation type and present their mutual coherence measurements as a box plot. 
    We used parameters $n = 256$, $\alpha = 1$, $d = 32 \times 32$.
    The V1 parameters were $s = 50$ and $f = 6$.
    }
    
    \label{fig:MCBox}
\end{figure}

\begin{figure}[h]
    \centering
    
    \includegraphics[width=\linewidth]{A_images/all_patches/png/DotProductHist_all_patches.png}
    
    \caption{
    Here we have the column-wise dot products, from $\mat{AA}^T$, viewed as a histogram.
    We used parameters $n = 256$, $\alpha = 1$, $d = 32 \times 32$.
    The V1 parameters were $s = 50$ and $f = 6$.
    }
    
    \label{fig:DotHist}
\end{figure}

\begin{figure}[h]
    \centering
    \includegraphics[width=\linewidth]{A_images/all_patches/png/SVD_256_patches.png}
    
    \caption{
    From the singular value decomposition of $\mat{A}$, we plotted the elements of $\bm{\Sigma}$ directly. We used parameters $n = 256$, $\alpha = 1$, $d = 32 \times 32$.
    The V1 parameters were $s = 50$ and $f = 6$.
    }
    \label{fig:SVD}
\end{figure}

\begin{figure}[h]
    \centering

   \includegraphics[width=\linewidth]{A_images/all_patches/png/all_patches_true_pc_per_method.png}
   
    \caption{Here is a scatter plot of the element of $\vec{p}$ for several patches of \textit{Barbara}.
    We believe the ``hills" evident in patches 169 and 233 cause the exorbitant error found in Fig.~\ref{fig:PCError}.
    We used parameters $n = 256$, $\alpha = 1$, $d = 32 \times 32$.
    The V1 parameters were $s = 50$ and $f = 6$.}
    \label{fig:PCScatter}
\end{figure}

\begin{figure}[h]
    \centering
    
    \includegraphics[width=\linewidth]{A_images/all_patches/png/all_patches_pc_scatter.png}
    
    \caption{
    We compare $\mat{p}$ to $\hat{\mat{p}}$ of several patches of \textit{Barbara}.
    We used parameters $n = 256$, $\alpha = 1$, $d = 32 \times 32$.
    The V1 parameters were $s = 50$ and $f = 6$.
    The dotted line represents $x=y$.
    }
    \label{fig:PCCompareScatter}
\end{figure}

\begin{figure}[h]
    \centering
    \includegraphics[width=\linewidth]{A_images/all_patches/png/all_patches_error_cumsum.png}
    
    \caption{We looked at the cumulative sum of errors between $\mat{p}$ from \textit{Barbara} and $\hat{\mat{p}}$ from the reconstructions.
    We used parameters $n = 256$, $\alpha = 1$, $d = 32 \times 32$.
    The V1 parameters were $s = 50$ and $f = 6$.
    }
    \label{fig:PCError}
\end{figure}

\subsection{Extra}

\subsubsection{Dimensionality}

To compute the dimensionality $D$ of $\mat{A}$, 
\begin{equation}
    D = \left( \sum \frac{\sigma}{\sum \sigma_i}\right)^{-1}
\end{equation}
where $\sigma_i$ is the singular value of column $i$ in $\bm\Sigma$.

\begin{figure}[h]
    \centering
   \includegraphics[width=\linewidth]{A_images/all_patches/png/Dimensions_all_patches.png}
    \caption{
    We found the dimensions $D$ from 100 different design matrices per observation type.
    We used parameters $n = 256$, $\alpha = 1$, $d = 32 \times 32$.
    The V1 parameters were $s = 50$ and $f = 6$.
    }
    \label{fig:Dimension}
\end{figure}

\begin{figure}[h]
    \centering
    \includegraphics[width=\linewidth]{A_images/all_patches/png/pc_top3_images_256_patch_58.png}
    
    \caption{We render $\hat{\mat{p}}_{58}$ as pixel images. 
    Each vector is reminiscent of it's observation type taken from \textit{Barbara.}
    We used parameters $n = 256$, $\alpha = 1$, $d = 32 \times 32$.
    The V1 parameters were $s = 50$ and $f = 6$.
    \todo{are we keeping the PC images in or not? could be a way to explain the "blobbyness" thing in the intro}}
    
    \label{fig:PCImages}
\end{figure}

\end{document}
